% ============================================================================
% CHAPTER 3: ANALYSIS OF EXISTING WORK AND LIMITATIONS
% ============================================================================
\chapter{Analysis of Existing Work and Limitations}

\section{Working of Existing System}

Existing AI-assisted research systems operate through a fundamentally sequential, human-in-the-loop paradigm. The typical workflow proceeds as follows:

\textbf{Step 1 --- Manual Query Formulation:} The researcher constructs search queries using domain-specific keywords and Boolean operators. Tools like Semantic Scholar accept natural language queries but rely on embedding-based retrieval (SPECTER, 768-dim), which achieves Recall@20 of 73.4\% on the TREC-COVID benchmark---meaning roughly 1 in 4 relevant papers is missed in the initial retrieval.

\textbf{Step 2 --- Individual Paper Review:} Each retrieved paper must be individually accessed, read, and analyzed. At an average reading speed of 3--5 pages per hour for technical literature, reviewing 50 papers requires 100--200 hours of focused effort.

\textbf{Step 3 --- Manual Synthesis:} The researcher manually identifies patterns, contradictions, and gaps across the reviewed literature. This synthesis step has no tooling support in existing systems---it relies entirely on human cognitive capacity.

\textbf{Step 4 --- Writing and Formatting:} The final literature review, methodology proposal, and report must be manually written, formatted, and checked for consistency. Citation management tools (Zotero, Mendeley) assist with reference formatting but provide no content generation.

Fig~\ref{fig:existing_workflow} illustrates this sequential workflow.

\begin{figure}[H]
\centering
\begin{tikzpicture}[
    node distance=1.5cm,
    block/.style={rectangle, draw, fill=blue!10, text width=3.5cm, minimum height=1.2cm, align=center, rounded corners, font=\small},
    arrow/.style={-{Stealth[length=3mm]}, thick}
]
\node[block] (q) {Manual Query\\Formulation};
\node[block, right=of q] (r) {Paper-by-Paper\\Review (100-200h)};
\node[block, right=of r] (s) {Manual\\Synthesis};
\node[block, right=of s] (w) {Report Writing\\and Formatting};
\draw[arrow] (q) -- (r);
\draw[arrow] (r) -- (s);
\draw[arrow] (s) -- (w);
\end{tikzpicture}
\caption{Existing Sequential Research Workflow}
\label{fig:existing_workflow}
\end{figure}

\section{Limitation of Existing System}

The limitations of existing systems are categorized into five dimensions:

\textbf{1. Single-Perspective Bias:} All existing tools produce outputs from a single model invocation. This introduces systematic bias toward the model's training data distribution and temperature-dependent sampling artifacts. Research by Wang et al.\ (2023) demonstrated that single-shot LLM outputs exhibit 18\% higher variance in quality compared to ensemble approaches, with a standard deviation of $\pm$12 points on a 100-point quality scale.

\textbf{2. No Quality Assurance:} No existing system implements automated quality scoring of generated outputs. Elicit provides ``confidence indicators'' based on claim extraction heuristics, but these measure extraction confidence (is this sentence a claim?), not content quality (is this claim accurate, specific, and well-supported?). The absence of self-evaluation means that low-quality outputs---vague summaries, hallucinated citations, generic methodology descriptions---are presented alongside high-quality outputs with no differentiation.

\textbf{3. Zero Visualization:} No existing research assistant generates system architecture diagrams, methodology flowcharts, or comparative visualization artifacts. Researchers must use separate tools (draw.io, Lucidchart, Mermaid, PlantUML) and manually translate research findings into visual formats---a process requiring 4--8 hours per comprehensive diagram.

\textbf{4. Fragmented Pipeline:} Current tools address isolated steps of the research workflow. A typical setup requires: Semantic Scholar for discovery + Zotero for citation management + Elicit for data extraction + ChatGPT for writing assistance + draw.io for diagrams + LaTeX/Word for formatting = 6 separate tools with no data integration between them.

\textbf{5. Connection Instability:} Web-based AI tools suffer from timeout issues during long-running operations. Standard HTTP requests timeout after 30--60 seconds, while comprehensive research synthesis requires 5--15 minutes of continuous processing. This forces users into fragmented, multi-session workflows that lose context between interactions.

\textbf{6. No Methodology Generation:} Existing systems can summarize what research exists but cannot propose how to address identified gaps. Generating a technically detailed methodology proposal---with specific architectures, loss functions, training protocols, and quantitative targets---requires domain expertise that goes beyond summarization.

Table~\ref{tab:limitations} quantifies these limitations against our proposed system.

\begin{table}[H]
\centering
\caption{Quantitative Comparison of Limitations}
\label{tab:limitations}
\begin{tabularx}{\textwidth}{|X|c|c|c|}
\hline
\textbf{Metric} & \textbf{Semantic Scholar} & \textbf{Elicit} & \textbf{Our System} \\
\hline
Time for 50-paper review & 100--200 hours & 8--15 hours & 10--15 minutes \\
\hline
Perspectives analyzed & 1 & 1 & 3 (parallel) \\
\hline
Output quality scoring & None & Heuristic & LLM-Judge (0--100) \\
\hline
Diagrams generated & 0 & 0 & 6 types, 8 themes \\
\hline
Report word count & 0 & 200--500 & 6000+ \\
\hline
Citations formatted & Manual & Partial & 30--40+ auto \\
\hline
Methodology proposals & 0 & 0 & 5--8 per topic \\
\hline
Max pipeline duration & N/A & 30s timeout & 15 min (stable) \\
\hline
\end{tabularx}
\end{table}
